\documentclass[11pt]{article}

\usepackage[margin=1in]{geometry}
\usepackage{amsmath, amssymb, amsthm}
\usepackage{bm}
\usepackage{hyperref}
\usepackage{xcolor}
\usepackage{booktabs}

\hypersetup{colorlinks=true, linkcolor=blue!50!black, urlcolor=blue!50!black}

\newcommand{\E}{\mathbb{E}}
\newcommand{\Eq}[1][q]{\mathbb{E}_{#1}}
\newcommand{\Eqx}{\mathbb{E}_{q(x_i)}}
\newcommand{\Eqxx}{\mathbb{E}_{q(x_i,\,x_{i+1})}}
\newcommand{\Cov}{\operatorname{Cov}}
\newcommand{\diag}{\operatorname{diag}}
\newcommand{\tr}{\operatorname{tr}}
\newcommand{\GP}{\operatorname{GP}}
\newcommand{\given}{\,|\,}

\newcommand{\code}[1]{\texttt{\small #1}}

\title{The EFGP-SING E-step without Monte Carlo\\
\large Closed form via a custom Gaussian-mixture spreader}
\author{}
\date{}

\begin{document}
\maketitle

\begin{abstract}
\noindent
The EFGP-SING E-step needs, for each output dim $r$, the marginal
$q(x)$-expectations that build the precision $A_r$ and right-hand side
$h_r$ (Eqs.~\eqref{eq:Ah}--\eqref{eq:hr-stein}).  This note makes one
observation: every such expectation is an instance of
$\E_q[\sum_i w_i\,e^{2\pi i\,x_i^\top\xi}]$, the expectation of a
nonuniform Fourier sum over the latent points, and \emph{because
expectation and the Fourier transform are both linear and commute},
this equals the Fourier transform of the \emph{expected} source density
$g(x)=\sum_i w_i\,\mathcal{N}(x;m_i,S_i)$.  The entire E-step is
therefore ``deposit $g$ on a grid and FFT it.''  No per-source
characteristic function is ever evaluated --- the FFT performs the
averaging; the Gaussian CF is merely the per-source \emph{description}
of the same target, useful for reading off structure (BTTB, the cheap
derivative rule) but never a computational ingredient.

This is realised as a custom Gaussian-mixture spreader, now the default
EFGP-SING E-step (\code{estep\_method='gmix'}), replacing the earlier
Monte-Carlo build (a pseudo-cloud $x_i^{(s)}\!\sim\!q(x_i)$ through a
single Type-1 NUFFT; see \texttt{efgp\_estep.tex}).  It is
deterministic (no MC noise into the M-step), preserves the BTTB
structure of $A_r$, and approximates the target by equispaced
quadrature --- its only errors are FFT discretisation and
Gaussian-stencil truncation, both controllable.  Cost is
$O(T \cdot (\sigma_{\max}/\sigma_{\min})^D + N^D \log N)$ per FFT-set
call.  Empirically, on a GP-drift recovery benchmark with fixed
emissions, it recovers $\ell$ and $\sigma^2$ slightly closer to truth
than the $S\!=\!2$ MC default, yields slightly lower drift error, and
--- unlike MC --- has no catastrophic-divergence failure mode at large
$T$ (\S\ref{sec:bench}).
\end{abstract}

\tableofcontents

% ============================================================================
\section{The expectations we need}
\label{sec:setup}
% ============================================================================

We adopt the notation of \texttt{efgp\_estep.tex} (Layers 1--3).  The
mean-field optimum for $q(w_r)=\mathcal{N}(\mu_r, A_r^{-1})$ is the
system
\begin{equation}
A_r\,\mu_r \;=\; h_r,
\label{eq:Ah-system}
\end{equation}
with
\begin{equation}
A_r \;=\; I + \frac{1}{\sigma_r^2}\sum_i \Delta_i\,
\Eqx\!\bigl[\phi(x_i)\phi(x_i)^{*}\bigr],
\qquad
h_r \;=\; \frac{1}{\sigma_r^2}\sum_i
\Eqxx\!\bigl[\phi(x_i)\,(x_{i+1,r}-x_{i,r})\bigr],
\label{eq:Ah}
\end{equation}
and Fourier features
\begin{equation}
\phi_\theta(x) \;=\; D_\theta\, F_x,
\qquad
F_x[k] \;=\; e^{2\pi i\, x^\top \xi_k}.
\end{equation}
Stein's lemma collapses the pairwise expectation in $h_r$ to marginal
expectations under $q(x_i)$ plus a covariance correction:
\begin{equation}
h_r \;=\; \frac{1}{\sigma_r^2}\sum_i
\Bigl\{\Eqx[\phi(x_i)]\,d_{i,r}
\;+\; \Eqx[J_\phi(x_i)]^{*}\,C_{i,\cdot r}\Bigr\},
\qquad
d_i := m_{i+1}-m_i,\ \ C_i := S_{i,i+1}-S_i.
\label{eq:hr-stein}
\end{equation}
Every quantity in \eqref{eq:Ah}--\eqref{eq:hr-stein} is a
\emph{marginal} expectation under $q(x_i)\!=\!\mathcal{N}(m_i, S_i)$ of
a complex exponential or its derivative.  These are all instances of
one primitive --- the expectation of a Fourier sum --- which
\S\ref{sec:cf} collapses to a single FFT of a density.

% ============================================================================
\section{The one primitive: $\E[\text{Fourier sum}]$ is an FFT of a density}
\label{sec:cf}
% ============================================================================

Collect the quantities of \S\ref{sec:setup} into a single object: each
is
\begin{equation}
F(\xi)\;=\;\Eq\Big[\sum_i w_i\,e^{2\pi i\,x_i^\top\xi}\Big],
\label{eq:primitive}
\end{equation}
for some weights $w_i$ and output frequencies $\xi$ on a regular grid
(the spectral grid $\{\xi_k\}$ for $h_r$, the difference grid
$\{\xi_k-\xi_\ell\}$ for $A_r$).  The bracketed sum is a nonuniform
Fourier sum over the random points $\{x_i\}$ --- a Type-1 NUFFT of point
masses --- and the whole E-step reduces to taking its \emph{expectation}.

\paragraph{Expectation commutes with the Fourier transform.}
Expectation is an integral and the Fourier transform is an integral, so
the two commute: $\E$ slides \emph{through} the transform and lands on
the source distribution rather than on each exponential.  Step by step
--- take the expectation inside the sum, write each
$\E[e^{2\pi i\,x_i^\top\xi}]=\int \mathcal{N}(x;m_i,S_i)\,e^{2\pi
i\,\xi^\top x}\,dx$ as the Fourier transform of the $i$-th density (an
integral, no closed form needed), and pull the sum back inside:
\begin{equation}
\boxed{\;
F(\xi)
\;=\;\sum_i w_i\,\E\big[e^{2\pi i\,x_i^\top\xi}\big]
\;=\;\int\!\Big(\underbrace{\textstyle\sum_i w_i\,\mathcal{N}(x;m_i,S_i)}
       _{=:\;g(x)}\Big) e^{2\pi i\,\xi^\top x}\,dx
\;=\;\widehat g(\xi).
\;}
\label{eq:commute}
\end{equation}
The density $g$ is the \emph{expected point cloud} --- each random point
replaced by its Gaussian smear:
\begin{equation}
g(x)\;=\;\Eq\Big[\sum_i w_i\,\delta(x-x_i)\Big]
\;=\;\sum_i w_i\,\mathcal{N}(x;\,m_i, S_i).
\label{eq:gmix-density}
\end{equation}
So the entire E-step primitive is one instruction: \emph{deposit the
expected density $g$ on a grid and FFT it}, $\widehat g(\xi)\approx
\mathrm{FFT}[g]$ on the dual grid (\S\ref{sec:gmix-spreader}).  Nothing
per-source is expanded in closed form; the FFT performs the averaging,
because it is linear.

\paragraph{Same target as Monte Carlo, two commuting operations in the
other order.}  This is also the exact relationship to the MC E-step
(\texttt{efgp\_estep.tex}) that \code{estep\_method='gmix'} replaces:
\begin{itemize}
\item \textbf{MC}: \emph{sample} the cloud, then FFT ---
$\tfrac1S\sum_s\mathrm{NUFFT}\big[\sum_i w_i\,\delta(\cdot-x_i^{(s)})\big]$.
\item \textbf{gmix}: take the \emph{expectation first} ($\to g$), then FFT.
\end{itemize}
They target the identical $F$ and agree precisely because $\E$ and FFT
commute; gmix does the averaging analytically by depositing $g$, so it
is deterministic --- no shot noise enters the M-step.

\paragraph{Where the characteristic function hides.}
The per-source factor $\Eq[e^{2\pi i\,x_i^\top\xi}]$ that
\eqref{eq:commute} left as an integral has the closed-form value
$e^{2\pi i\,m^\top\xi}\,e^{-2\pi^2\,\xi^\top S\,\xi}$ --- the Gaussian
characteristic function.  Substituting it (and reattaching the feature
normaliser $D_\theta$, using $\partial_{x_j}\phi_k = 2\pi
i\,\xi_{k,j}\,\phi_k$) gives the equivalent \emph{spectral-side} forms
\begin{align}
\Eqx[\phi_k(x)]
&\;=\; D_{\theta,k}\,e^{2\pi i\,m^\top \xi_k}\,e^{-2\pi^2\,\xi_k^\top S\,\xi_k},
\label{eq:Ephi}
\\[2pt]
\Eqx[\phi_k(x)\,\overline{\phi_\ell(x)}]
&\;=\; D_{\theta,k}\,\overline{D_{\theta,\ell}}\,
e^{2\pi i\, m^\top(\xi_k-\xi_\ell)}\,
e^{-2\pi^2\,(\xi_k-\xi_\ell)^\top S\,(\xi_k-\xi_\ell)},
\label{eq:Ephiphi}
\\[2pt]
\Eqx\!\bigl[\partial_{x_j}\phi_k(x)\bigr]
&\;=\; 2\pi i\,\xi_{k,j}\,\Eqx[\phi_k(x)].
\label{eq:EJphi}
\end{align}
These are a useful \emph{description} of the target $F$ --- they expose
the per-source damping $e^{-2\pi^2\xi^\top S_i\xi}$, the BTTB structure
(\S\ref{sec:Ar}), and the cheap derivative rule \eqref{eq:EJphi} ---
but they are \emph{never evaluated}.  The algorithm computes $F$ as
$\widehat g$ by one FFT; the damping factor \emph{emerges} as the
transform of the deposited density, not from its closed form.  The only
closed form the code consumes is the Gaussian \emph{density} in
\eqref{eq:gmix-density} (the spread weights), the other half of the
same Fourier pair.

Two structural facts the spectral-side forms make visible.  First,
\eqref{eq:Ephiphi} is \emph{not} the product of \eqref{eq:Ephi} with
its conjugate --- the damping sees $\xi_k - \xi_\ell$, not $\xi_k$ and
$\xi_\ell$ separately.  This is the single most important consequence
of $x$ being uncertain, and it is what makes $A_r$ BTTB in the
frequency \emph{differences} (\S\ref{sec:Ar}).  Second, the derivative
rule \eqref{eq:EJphi} means the Stein-correction term in $h_r$ is the
\emph{same} $F$ multiplied by $2\pi i\,\xi_j$ on the spectral grid: one
extra coefficient grid through the same FFT, with no new density and no
per-source work.

% ============================================================================
\section{The precision $A_r$: BTTB is preserved}
\label{sec:Ar}
% ============================================================================

From \eqref{eq:Ah}, the precision is built from
$\Eqx[\phi_k(x)\overline{\phi_\ell(x)}]$, which is the expectation of a
\emph{single} exponential at the difference frequency:
\begin{equation}
\Eqx[\phi_k\overline{\phi_\ell}]
\;=\; D_{\theta,k}\,\overline{D_{\theta,\ell}}\,
\Eqx\!\big[e^{2\pi i\,x^\top(\xi_k-\xi_\ell)}\big].
\end{equation}
So it depends on $(k,\ell)$ only through $\Delta\xi:=\xi_k-\xi_\ell$ ---
for \emph{any} $q$, before any Gaussian assumption.  Hence, on an
equispaced spectral grid, $A_r = I + D_\theta T_r D_\theta$ is BTTB,
exactly as in the MC formulation, with generator
\begin{equation}
\boxed{\;
T_r[\Delta\xi]
\;=\;
\sum_i \frac{\Delta_i}{\sigma_r^2}\,
\Eqx\!\big[e^{2\pi i\,x_i^\top \Delta\xi}\big]
\;=\;
\widehat{g_T}(\Delta\xi),
\qquad
g_T := \sum_i \frac{\Delta_i}{\sigma_r^2}\,\mathcal{N}(\cdot;m_i,S_i).
\;}
\label{eq:Tr-generator}
\end{equation}
This is exactly the \S\ref{sec:cf} primitive \eqref{eq:commute} with
weights $w_i=\Delta_i/\sigma_r^2$, evaluated on the difference grid:
the FFT of one deposited density $g_T$.  (The per-source expectation
expands to the spectral-side CF $\sum_i (\Delta_i/\sigma_r^2)\,
e^{2\pi i\,m_i^\top\Delta\xi}\,e^{-2\pi^2\,\Delta\xi^\top S_i\,\Delta\xi}$,
but --- as in \S\ref{sec:cf} --- that form is never evaluated.)  Matvecs
$A_r v$ remain $O(M\log M)$ via FFT on the generator; nothing about
\code{make\_A\_apply} or the CG solver changes.

\paragraph{Why it is not a NUFFT of the means.}
The MC version forms $T_r$ from a single Type-1 NUFFT of a sampled
point cloud, $O(N + M\log M)$.  The generator \eqref{eq:Tr-generator}
is \emph{not} a NUFFT of the means $\{m_i\}$: the sources are not points
but the smeared density $g_T$, and the per-source spread of
$\mathcal{N}(\cdot;m_i,S_i)$ --- equivalently the damping
$e^{-2\pi^2\,\Delta\xi^\top S_i\,\Delta\xi}$ --- couples the output
frequency $\Delta\xi$ and the source $i$ together.  Pulling it outside a
single NUFFT of the means would require $S_i$ independent of $i$.  The
resolution is precisely \S\ref{sec:cf}: it is the FFT of the expected
density, carried out by the spreader of \S\ref{sec:gmix-spreader}.

% ============================================================================
\section{The right-hand side $h_r$}
\label{sec:hr}
% ============================================================================

The same primitive builds $h_r$, now on the spectral grid $\{\xi_k\}$.
The mean-increment term of \eqref{eq:hr-stein} uses $\Eqx[\phi_k]$:
\begin{equation}
(h_{1,r})_k \;=\; \frac{D_{\theta,k}}{\sigma_r^2}\sum_i d_{i,r}\,
\Eqx\!\big[e^{2\pi i\,x_i^\top\xi_k}\big]
\;=\; \frac{D_{\theta,k}}{\sigma_r^2}\,\widehat{g_{h_1}}(\xi_k),
\qquad
g_{h_1} := \sum_i d_{i,r}\,\mathcal{N}(\cdot;m_i,S_i).
\label{eq:h1-cf}
\end{equation}
The Stein correction enters through the adjoint Jacobian.  By the
derivative rule \eqref{eq:EJphi}, $\partial_{x_j}$ acts as the
multiplier $2\pi i\,\xi_{k,j}$ on the spectral grid (conjugating to form
$J_\phi^{*}$ flips the sign), so $h_{2,r}$ is the \emph{same} FFT output
weighted by $C$ and post-multiplied by $-2\pi i\,\xi_j$:
\begin{equation}
(h_{2,r})_k
\;=\; \frac{D_{\theta,k}}{\sigma_r^2}\sum_{j=1}^{D_x}
(-2\pi i\,\xi_{k,j})\,\widehat{g_{h_2,j}}(\xi_k),
\qquad
g_{h_2,j} := \sum_i [C_i]_{j,r}\,\mathcal{N}(\cdot;m_i,S_i).
\label{eq:h2-cf}
\end{equation}
Both are the \S\ref{sec:cf} primitive on the spectral grid --- the same
object as the generator \eqref{eq:Tr-generator}, just at $\{\xi_k\}$
instead of $\{\Delta\xi\}$, with weights $d_{i,r}$ (resp.\ $[C_i]_{j,r}$)
in place of $\Delta_i/\sigma_r^2$.  Each is one deposited density
through one FFT; the Jacobian costs only the extra $2\pi i\,\xi_j$
multiplier, not a new spread.  Setting $S_i\!\to\!0$ collapses each
$\mathcal{N}$ to a $\delta$ and recovers the MC build exactly.

% ============================================================================
\section{How much of this is Fourier-specific?}
\label{sec:fourier-specific}
% ============================================================================

Before the algorithm, it is worth separating what the construction owes
to the \emph{features} being Fourier from what it owes to nothing in
particular.  The split is clean and pinpoints exactly what the basis
buys.

\paragraph{The smearing is basis-free.}
The reframing of \S\ref{sec:cf} --- that each expected sufficient
statistic is a functional of the \emph{expected} density
$g=\sum_i w_i\,\mathcal{N}(\cdot;m_i,S_i)$ --- is pure linearity of
expectation and uses nothing about Fourier:
\begin{equation}
\Eq\Big[\sum_i w_i\,\phi_k(x_i)\Big]
\;=\; \sum_i w_i\!\int \phi_k(x)\,\mathcal{N}(x;m_i,S_i)\,dx
\;=\; \langle \phi_k,\,g\rangle,
\label{eq:basis-free}
\end{equation}
for \emph{any} features $\phi_k$.  ``Smear each uncertain point by its
posterior covariance, then evaluate features against the smeared
cloud'' would read identically for random, kernel, or neural features.
Even the per-source closed form is not special to Fourier: RBF,
arccos/ReLU, and polynomial features all integrate against a Gaussian
in closed form --- the Gaussian characteristic function is merely the
instance for $\phi_k=e^{2\pi i\,x^\top\xi_k}$.

\paragraph{But the smearing is not the sufficient statistics.}
The smear produces $g$ in \emph{input} space.  The expected sufficient
statistics $A_r,h_r$ live in \emph{feature} space, and in the Fourier
basis they \emph{are} the Fourier coefficients of $g$:
$h_r=\widehat{g_{h_1}}(\xi_k)$ on the spectral grid \eqref{eq:h1-cf},
and $A_r$'s generator $=\widehat{g_T}(\xi_k-\xi_\ell)$ on the difference
grid \eqref{eq:Tr-generator}.  So forming the statistics is \emph{not}
the smear; it is the FFT that carries $g$ from input space to feature
space.  The smear is the FFT's \emph{input}; the statistics are its
\emph{output} --- different objects related by the transform.  (One
could instead form them by direct summation of \eqref{eq:Ephi},
$O(MT)$; the FFT is the fast evaluator of ``Fourier coefficients of
gridded data,'' which is what the statistics are.)

\paragraph{Three roles, one basis-free.}
The pipeline has three steps, and only the first survives a change of
features:
\begin{center}
\renewcommand{\arraystretch}{1.3}
\begin{tabular}{lll}
\toprule
Step & Produces & Fourier-specific? \\
\midrule
Smear & pseudo-obs density $g$ (input space) & \textbf{no} (linearity) \\
FFT of $g$ & suff.\ stats $A_r$-generator, $h_r$ (feature space) & \textbf{yes} \\
Solve $A_r\mu_r=h_r$ & the $q(f)$ posterior & \textbf{yes} (BTTB) \\
\bottomrule
\end{tabular}
\end{center}
The FFT step and the BTTB structure of $A_r$ both rest on the
\emph{character} property $\phi_k\,\overline{\phi_\ell}=\phi_{k-\ell}$:
products of Fourier features are again Fourier features, indexed by the
\emph{difference}.  This is why $\Eq[\phi_k\overline{\phi_\ell}]=
\widehat g(\xi_k-\xi_\ell)$ depends on $(k,\ell)$ only through the
difference (\S\ref{sec:Ar}), and it is strictly stronger than ``a fast
transform exists'': even translated-kernel features on a grid
($\phi_k(x)=\kappa(x-z_k)$, which \emph{do} admit an FFT) fail it,
because $\Eq[\kappa(\cdot-z_k)\,\kappa(\cdot-z_\ell)]$ depends on both
indices once $q$ is localised.

\paragraph{Given the pseudo-observations, the Toeplitz structure is the
headline.}
The Toeplitz property does double duty, and the two uses are
inseparable.  It is what lets $A_r$ be \emph{represented} by an $O(M)$
generator (rather than a dense $O(M^2)$ Gram one would have to
materialise) --- and that generator \emph{is} the sufficient statistic
$\widehat{g_T}$, filled in by one FFT --- \emph{and} it is what lets
each CG matvec $A_r v$ be an $O(M\log M)$ FFT.  ``Form the statistic''
and ``solve fast'' are the same Fourier coin.  Cost-wise the solve
dominates: forming the statistics is a fixed handful of FFTs per inner
iteration ($1+D_{\text{out}}(1+D_{\text{lat}})$, \S\ref{sec:impl}),
while the solve is $k_{\text{cg}}$ of them.  So, given the smeared
pseudo-observations, what the Fourier basis buys is the
$O(M)$-storage, $O(M\log M)$-matvec precision.

\paragraph{The clean control: same pseudo-obs, no FFT.}
SparseGP-SING (\code{sing/sing.py}) runs the \emph{same} E-step on the
\emph{same} smeared pseudo-observations, but in an inducing-point
basis: its sufficient statistics are dense Gram blocks formed by direct
kernel evaluation (no FFT) and solved by dense/low-rank algebra (no
Toeplitz).  Identical \eqref{eq:basis-free}; entirely different cost.
That neither the FFT-formation nor the Toeplitz-solve comes from the
smearing --- both are purchased by choosing Fourier features --- is
exactly what the EFGP-vs-SparseGP split demonstrates.

% ============================================================================
\section{Algorithm: per-source Gaussian-mixture spreader}
\label{sec:gmix-spreader}
% ============================================================================

The objects we need (\S\S\ref{sec:Ar}--\ref{sec:hr}) are all the
\S\ref{sec:cf} primitive \eqref{eq:commute},
\begin{equation}
F(\xi)
\;=\;
\widehat g(\xi)
\;=\;
\int g(x)\,e^{2\pi i\,\xi^\top x}\,dx,
\qquad
g(x) \;:=\; \sum_i w_i\,\mathcal{N}(x;\,m_i, S_i),
\label{eq:gmix-target}
\end{equation}
the continuous Fourier transform of a Gaussian \emph{mixture} density
$g$ in physical space (equivalently the per-source CF sum $\sum_i w_i\,
e^{2\pi i\,m_i^\top\xi}e^{-2\pi^2\xi^\top S_i\xi}$, which we never
evaluate).  This section is how the ``$\approx\mathrm{FFT}$'' of
\eqref{eq:commute} is actually carried out: build $g$ on a grid, FFT it.

\subsection{Spread the densities, FFT, done}
\label{sec:plain-story}

The algorithm is an equispaced quadrature of \eqref{eq:gmix-target}
followed by an FFT.  On a regular grid $\{x_a\}$ with spacing
$h_{\text{grid}}$ wide enough to enclose $g$, the Riemann sum
\begin{equation}
F(\xi) \;\approx\;
h_{\text{grid}}^{D}\sum_a g(x_a)\,e^{2\pi i\,\xi^\top x_a}
\label{eq:riemann}
\end{equation}
is, for output frequencies on the dual spectral grid (spacing
$h_{\text{spec}} = 1/(N\,h_{\text{grid}})$), exactly a $D$-dimensional
FFT of $g_a := g(x_a)$ up to fixed phase shifts.  So the only work is
building $g_a$.

\paragraph{Building $g_a$ is cheap.}
$\mathcal{N}(x_a;m_i,S_i)$ is negligible past $\sim n_\sigma\sigma_i$
from $m_i$, so each source contributes only to a stencil of radius
$r\approx n_\sigma\sigma_{\max}/h_{\text{grid}}$:
\begin{equation}
g_a \;=\; \sum_{i\,:\,a\in\text{stencil}(i)} w_i\,\mathcal{N}(x_a;m_i,S_i),
\end{equation}
costing $\sim T(2r+1)^D$ Gaussian evaluations rather than $T N^D$.  So
$g_a$ is a sparse stencil accumulator; one FFT then returns $F(\xi_k)$.
The spectral grid ($M_{\text{per dim}}^D$) and the BTTB difference grid
($(2M_{\text{per dim}}-1)^D$) are both centred crops of the same FFT
output, so one spread+FFT serves both.

\paragraph{It is a NUFFT with the deconvolution removed.}
A standard Type-1 NUFFT spreads each point mass with a \emph{fixed
artificial} kernel $\kappa$ (width set by target precision $\epsilon$),
FFTs, then \emph{deconvolves} by $\hat\kappa$ to undo the spreading.
Here the integrand is already a sum of smooth Gaussians, so we spread
the \emph{actual} per-source PDFs --- one line per stencil cell,
$g[a]\mathrel{+}=w_i\,\mathcal{N}(x_a;m_i,S_i)$ --- and \emph{skip the
deconvolution}: the resulting envelope
$\hat\kappa_i(\xi)=e^{-2\pi^2\xi^\top S_i\xi}$ is exactly the per-source
damping we wanted (the Gaussian CF of \S\ref{sec:cf}), not an artefact
to remove.  Outputs are read directly off the FFT grid, since EFGP's
frequencies are themselves gridded.  The price for using the true
$S_i$: the stencil radius tracks the data's $\sigma_{\max}$ instead of
a fixed $\epsilon$, so heterogeneous $\{S_i\}$ cost more cells
(\S\ref{sec:cost}).

\paragraph{Two errors.}
Quadrature (Riemann vs.\ integral; controlled by
$h_{\text{grid}}\lesssim\sigma_{\min}$ and a grid wide enough to hold
$g$) and stencil truncation (dropping cells past $r$; error
$\sim e^{-n_\sigma^2/2}$, quantified next).

\subsection{Stencil truncation error}
\label{sec:trunc}

Outside the stencil the Gaussian decays as
$\exp(-r^2 h_{\text{grid}}^2 / (2 \sigma_{\max}^2))$.  Truncating at
$r = 4\sigma_{\max}/h_{\text{grid}}$ gives relative error $\sim
e^{-8} \approx 3\times 10^{-4}$; matching standard NUFFT precision
$10^{-7}$ requires $r \approx 5.7\sigma_{\max}/h_{\text{grid}}$.  In
the implementation, $r$ is set conservatively from the largest
$\lambda_{\max}(S_i)$ across $i$.

% ============================================================================
\section{Cost analysis}
\label{sec:cost}
% ============================================================================

Per spread call, the work decomposes as
\begin{equation}
\boxed{\;
\text{spread cost}\ =\ T \cdot (2r+1)^D
\;+\;
N^D \log N.
\;}
\end{equation}
The FFT term is $O(N^D \log N)$ as in any NUFFT.  The spread term is
where the cost picture differs from a standard NUFFT.

\subsection{The fundamental cost driver}

The fine-grid spacing $h_{\text{grid}}$ is constrained by two requirements:

\begin{enumerate}
\item \textbf{Resolution} of the sharpest source: $h_{\text{grid}}
\lesssim \sigma_{\min}$ (otherwise the spread under-samples the
sharpest Gaussians, which then alias as constants in the FFT).
\item \textbf{Spatial coverage}: $N\,h_{\text{grid}} \geq$ source
extent $+ \sim 4 \sigma_{\max}$ (otherwise wide-tail Gaussians wrap
around the periodic FFT boundary).
\end{enumerate}

Picking $h_{\text{grid}} \sim \sigma_{\min}$ to satisfy~(1), the
stencil radius becomes
\begin{equation}
r \;\sim\; \frac{4\sigma_{\max}}{h_{\text{grid}}}
\;\sim\; 4 \cdot \frac{\sigma_{\max}}{\sigma_{\min}}.
\end{equation}
Per-source cell count $(2r+1)^D \sim (8\sigma_{\max}/\sigma_{\min})^D$.

\paragraph{Comparison to standard NUFFT.}
A standard NUFFT chooses its kernel width $w$ purely from the desired
precision $\epsilon$, independent of data: $(2w+1)^D$ cells per source
with $w \sim 7$.  Our spreader instead inherits the per-source width
from the integrand:
\begin{equation}
\frac{\text{GMIX cost}}{\text{NUFFT cost}}
\;\sim\;
\left(\frac{\sigma_{\max}/\sigma_{\min}}{w/4}\right)^D.
\end{equation}
For $D\!=\!2$ and $w\!=\!7$: GMIX matches NUFFT cost when
$\sigma_{\max}/\sigma_{\min} \approx 1.75$.  When the variance ratio
across sources is larger, GMIX pays a $(\sigma_{\max}/\sigma_{\min})^2$
factor over a standard NUFFT.

\paragraph{Implication.}
GMIX is exact (variant~A's only approximation is stencil truncation),
but its cost depends on the heterogeneity of $\{S_i\}$ across $i$.  In
single-scale regimes ($S_i$ approximately equal), it matches a
standard NUFFT.  In multi-scale regimes (sharp posteriors at some $t$,
broad at others), it pays a constant-factor cost over MC.

\subsection{Optimal stencil radius: $r\!\approx\!8$, fixed}
\label{sec:r-empirics}

Wall time is \emph{non-monotone} in $r$: $r\!=\!4$ runs the EM in 50 s,
$r\!=\!8$ in 36 s, $r\!=\!16$ in 53 s.  Smaller $r$ coarsens the BTTB
$T_r$, inflating the CG iteration count; the optimum balances stencil
cost $(2r+1)^D$ against CG work, landing at $r\!=\!8$
($n_\sigma\!\approx\!1.5$) for $D\!=\!2$.  Adapting $r$ as
$\sigma_{\max}$ shrinks during EM does not pay --- $r$ is a static JIT
argument, so each change forces a $\sim$10--15 s recompile, more than
the $\sim$5 s it saves.  \textbf{Default:} pick $r$ once from $V_0$ at
$n_\sigma\!=\!1.5$, floor at 8, freeze; \code{gmix\_adapt\_stencil=True}
opts into the adaptive variant.

% ============================================================================
\section{Implementation}
\label{sec:impl}
% ============================================================================

\subsection{New module: \code{sing/efgp\_gmix\_spreader.py}}

Pure-JAX primitives:
\begin{itemize}
\item \code{\_spread\_2d(m, S, w, xcen, h\_grid, N, r)} --- vmap'd
per-source stencil evaluation + scatter-add into an $N{\times}N$
complex grid.  Anisotropic Gaussian (uses the full $D{\times}D$
$S_i$, including off-diagonals).
\item \code{gmix\_nufft\_2d(...)} --- spread + \code{ifftshift /
ifft2 / fftshift} + Riemann-sum normalisation $(N h_{\text{grid}})^D
= (1/h_{\text{spec}})^D$ + centering phase $e^{2\pi i \xi^\top
xcen}$.
\item \code{stencil\_radius\_for(S, h\_grid, n\_sigma=4.0)} --- picks
$r$ from $\max_i \lambda_{\max}(S_i)$.
\item \code{pick\_grid\_size(h\_spec, m\_extent, sigma\_max)} --- picks
$N$ as the next power of 2 satisfying both resolution and coverage
constraints (\S\ref{sec:cost}).
\end{itemize}

\subsection{New q(f) update: \code{compute\_mu\_r\_gmix\_jax}}

In \code{sing/efgp\_jax\_drift.py}, the deterministic (spread+FFT)
analogue of the MC \code{compute\_mu\_r\_jax}:
\begin{itemize}
\item Builds the BTTB generator from \eqref{eq:Tr-generator} via one
spread+FFT, cropped to the difference grid.
\item Builds $h_{1,r}$ and $h_{2,r}$ per output dim $r$ via additional
spread+FFTs cropped to the spectral grid.
\item Conjugates to align with the existing MC code's iflag=$-1$
NUFFT convention (so the BTTB matrix and CG solver are unchanged).
\item Falls through to the existing \code{cg\_solve} and the rest of
the pipeline.
\end{itemize}
Total spread+FFT calls per inner iter: $1 + D_{\text{out}}(1 +
D_{\text{lat}})$ ($= 7$ for $D_{\text{lat}} = D_{\text{out}} = 2$).

\subsection{Driver hook: \code{estep\_method}}

\code{fit\_efgp\_sing\_jax} now takes
\begin{itemize}
\item \code{estep\_method} $\in \{\code{'mc'}, \code{'gmix'}\}$,
default \code{'gmix'};
\item \code{gmix\_fine\_N}, \code{gmix\_stencil\_r} (auto-picked from
$V_0$ and $h_{\text{spec}}$ if omitted, then adapted once after
kernel warmup).
\end{itemize}
All three internal call sites (jit'd inner E-step scan, frozen-qf
path, M-step $\mu$-refresh) dispatch on this flag.

\subsection{Tests}

\begin{itemize}
\item \code{tests/test\_gmix\_spreader.py} --- spreader vs dense
closed-form reference on constant-$S$, varying-$S$,
resolution-convergence, and high-$S$-variation problems.  Errors
$10^{-7}$ to $10^{-4}$ depending on stencil truncation.
\item \code{tests/test\_efgp\_gmix\_vs\_mc.py} --- the deterministic
\code{compute\_mu\_r\_gmix\_jax} agrees with high-$S$-marginal MC
reference (8 seeds $\times$ $S\!=\!32$) to $\sim$1\,\%, while
single-seed $S\!=\!2$ MC sits at $\sim$12\,\%.  GMIX is $\sim 12\times$
closer to the MC reference than single-seed MC at production
settings, confirming it is the unbiased target the MC is
estimating.
\item \code{tests/test\_gmix\_nufft.py} --- earlier bucketed-NUFFT
correctness/complexity tests (variant~B) retained for reference.
\end{itemize}

% ============================================================================
\section{End-to-end benchmark}
\label{sec:bench}
% ============================================================================

\subsection{Setup}

\code{demos/bench\_gpdrift\_mc\_vs\_gmix.py}: GP-drift recovery in 2D,
$T \in \{500, 2000, 20000\}$, $\ell_{\text{true}} = 0.8$,
$\sigma^2_{\text{true}} = 1$, $\sigma_{\text{diff}} = 0.3$.  Drift
sampled fresh from an SE-GP with truth hypers; restoring term $-\alpha
x$ added for SDE stability.  Gaussian observations through a fixed $C
\in \mathbb{R}^{8\times 2}$, $R = 0.05 I$; emissions held fixed
(\code{learn\_emissions=False}) to isolate the q(f) regression block.
Both methods learn kernel hypers from $\ell_{\text{init}} = 0.3$
($0.4\times$ truth), $\sigma^2_{\text{init}} = 1$.  $n_\text{em} = 20{-}25$
outer iters, $n_\text{estep} = 10$ inner iters.  Default GMIX
configuration: $n_\sigma = 1.5$ (stencil-truncation eps $\sim e^{-1.1}
\approx 0.33$, ${\sim}\!2\times$ tighter than MC's $S\!=\!2$ shot noise);
no mid-run stencil adaptation; $r$ floored at 8.

\subsection{Results --- small/medium $T$}

At $T \in \{500, 2000\}$ the methods are close in both quality and
wall time; differences are small but consistent in GMIX's favour.

\begin{center}
\renewcommand{\arraystretch}{1.2}
\begin{tabular}{lrrrrrrr}
\toprule
$T$ & method & wall (s) & $\ell/\ell_\text{true}$ & $\sigma^2/\sigma^2_\text{true}$
& drift\_traj & drift\_grid & zero-baseline\\
\midrule
 500 & MC    &  30.9 & 1.06 & 0.76 & 0.148 & 0.353 & 0.676\\
 500 & GMIX  &  29.9 & 1.04 & 0.71 & 0.149 & 0.346 &      \\
2000 & MC    &  32.9 & 1.20 & 1.06 & 0.116 & 0.374 & 0.716\\
2000 & GMIX  &  39.9 & 1.14 & 0.97 & 0.111 & 0.350 &      \\
\bottomrule
\end{tabular}
\end{center}

\subsection{Results --- $T = 20{,}000$}

At $T = 20\text{k}$ the picture splits across seeds.  We ran two seeds
(controlling the GP-drift draw, the SDE simulation, and the
observation noise) plus an MC sample-count sweep.  At one seed, MC
catastrophically diverges; at another, MC works fine.  GMIX is
robust on both.

\begin{center}
\renewcommand{\arraystretch}{1.2}
\begin{tabular}{lcccrrr}
\toprule
seed & method & wall (s) & $\ell/\ell_\text{true}$ &
$\sigma^2/\sigma^2_\text{true}$ & drift\_traj & drift\_grid \\
\midrule
1 & MC, $S\!=\!2$  & 109 & 1.12 & 37.81 & \textbf{5.81} & \textbf{6.01} \\
1 & MC, $S\!=\!8$  & 100 & 0.99 & 31.03 & 4.67 & 4.69 \\
1 & MC, $S\!=\!16$ & 118 & 1.02 & 28.35 & 4.43 & 4.58 \\
1 & GMIX           & 124 & 1.06 &  0.71 & \textbf{0.044} & \textbf{0.226} \\
\addlinespace
2 & MC, $S\!=\!2$  & 100 & 1.19 &  0.87 & 0.033 & 0.346 \\
2 & MC, $S\!=\!8$  & 109 & 1.30 &  1.17 & 0.034 & 0.395 \\
2 & MC, $S\!=\!16$ & 105 & 1.29 &  1.15 & 0.033 & 0.390 \\
2 & GMIX           & 114 & 1.32 &  1.05 & 0.031 & 0.391 \\
\bottomrule
\end{tabular}
\end{center}

\paragraph{Two findings.}

\begin{enumerate}
\item \textbf{MC blow-up at $T\!=\!20$k is not deterministic but is real.}
On seed 1 MC's noise compounds through the natural-gradient updates,
$\sigma^2$ inflates 30--40$\times$, the kernel becomes too aggressive,
and the resulting drift posterior is meaningless (drift\_traj $\sim
5$, vs zero-baseline $0.77$).  On seed 2 the same setup produces a
clean fit.  We saw the failure on roughly half of attempts at this
problem size; smaller $T$ values do not exhibit it.
\item \textbf{More MC samples do not rescue the failure.}  Going from
$S\!=\!2$ to $S\!=\!8$ to $S\!=\!16$ (8$\times$ more compute on the
sample dimension) on seed 1 nudges drift\_traj from 5.81 to 4.43 ---
still two orders of magnitude worse than GMIX's 0.044.  $\sigma^2$
remains 28--38$\times$ truth.  The shot noise compounds through the
EM dynamics in a way that is not statistical (more samples per call)
but structural (each EM iter accumulates bias from the previous's
noisy update).
\end{enumerate}

GMIX has no analogous failure mode: its per-call output is
deterministic, so EM iter $k+1$ sees a clean response to iter $k$'s
update.  Both seeds converge GMIX to within $7\%$ of $\sigma^2_{\text{true}}$
and produce $1\!-\!\text{drift\_traj}/\text{zero}$-baseline $\geq$ 95\%.

\paragraph{Wall time at $T = 20$k.}  Within $\sim$15\% of MC across
seeds (114--124 s for GMIX vs 100--118 s for MC).  The slowdown
narrows from $\sim$1.5$\times$ at $T\!=\!2$k to ${\sim}1.1\times$ at
$T\!=\!20$k --- the per-iter overhead amortises better at larger $T$.

\subsection{Hyperparameter learning paths and the log-marginal landscape}

The hyperparameter paths plotted over the collapsed
log-marginal-likelihood landscape (computed under GMIX's converged
$q(x)$) appear in
\code{demos/\_bench\_mc\_vs\_gmix\_out/landscape\_paths\_T2000\_ls0.8.png}.
Both paths walk down the same valley to the same neighbourhood; the
GMIX trajectory is smoother (no MC noise) and lands slightly closer to
the true optimum.  At $T\!=\!20$k seed 1, MC's path drifts upward
toward $\sigma^2 \to \infty$ --- visible directly in the path plot.

% ============================================================================
\section{Where this sits among the alternatives}
\label{sec:alternatives}
% ============================================================================

For completeness, a comparison of the closed-form approximations
considered earlier:

\begin{center}
\renewcommand{\arraystretch}{1.4}
\begin{tabular}{lll}
\toprule
Variant & Per-spread cost & Error \\
\midrule
\textbf{(A) Custom spreader} (this note)
& $O(T(\sigma_{\max}/\sigma_{\min})^D + N^D \log N)$
& exact (FFT + stencil truncation) \\
(B) Bucketed NUFFT
& $O(BT/B + B M^D \log M) = O(T + B M^D \log M)$
& within-bucket bias from $S$-spread \\
(C) Gauss--Hermite ($K^D$ nodes)
& $O(TK^D + M^D \log M)$
& deterministic, $\sim (\xi^\top S\xi)^K$ at $K$ nodes \\
MC ($S$ samples)
& $O(TS + M^D \log M)$
& stochastic, $\sigma\sim\sqrt{1/S}$ \\
\bottomrule
\end{tabular}
\end{center}

Each row trades a different cost driver for a different error
mechanism.  Variant A removes both stochasticity and bucket bias at
the price of $(\sigma_{\max}/\sigma_{\min})^D$ in spread cost; MC
keeps the spread cheap but has shot noise that does not vanish at
fixed $S$.

% ============================================================================
\section{Recommendation}
\label{sec:reco}
% ============================================================================

\begin{itemize}
\item \textbf{Default to GMIX}.  It is exact (up to controllable FFT
truncation), gives slightly better drift error and hyperparameter
recovery, and has no MC noise to interact with the M-step.
\item \textbf{Use MC if cost matters more than precision.}  At $T \gg
1000$ in regimes where $\sigma_{\max}/\sigma_{\min}$ is large, MC
remains $\sim 1.5\times$ faster.  The MC noise at $S\!=\!2$ does
land within the same log-marginal valley; the differences are
constant-factor.
\item \textbf{Tune \code{kernel\_warmup\_iters}} so the single GMIX
stencil-radius adaptation lands at a sensible $\sigma_{\max}$.  Too
early and the adaptation triggers before the smoother sharpens;
too late and the cold-start runs all use a wide stencil.
\end{itemize}

% ============================================================================
\section{Math-to-code dictionary}
\label{sec:dict}
% ============================================================================

\noindent
\begin{tabular}{ll}
\toprule
Math & Code \\
\midrule
Primitive $F=\widehat g$, $\E[\phi]$ \eqref{eq:commute}, \eqref{eq:Ephi}
& \code{gmix\_nufft\_2d} \\
Generator $T_r=\widehat{g_T}$ \eqref{eq:Tr-generator}
& \code{gmix\_nufft\_2d} cropped to BTTB diff grid \\
$h_{1,r}=\widehat{g_{h_1}},\ h_{2,r}$ \eqref{eq:h1-cf}, \eqref{eq:h2-cf}
& \code{compute\_mu\_r\_gmix\_jax} \\
Per-source spread $\sum_i w_i \mathcal{N}(\cdot; m_i, S_i)$
& \code{\_spread\_2d} \\
Stencil radius from $\sigma_{\max}$
& \code{stencil\_radius\_for} \\
Fine grid size $N$ from $h_{\text{spec}}$ + $\sigma$
& \code{pick\_grid\_size} \\
Driver flag
& \code{fit\_efgp\_sing\_jax(estep\_method=...)} \\
Adaptive-$r$ once after warmup
& \code{\_maybe\_adapt\_stencil} in \code{efgp\_em.py} \\
Benchmark
& \code{demos/bench\_gpdrift\_mc\_vs\_gmix.py} \\
Landscape plot
& \code{demos/plot\_mc\_vs\_gmix\_landscape.py} \\
Tests
& \code{tests/test\_gmix\_spreader.py},\\
& \code{tests/test\_efgp\_gmix\_vs\_mc.py} \\
\bottomrule
\end{tabular}

\medskip
\noindent
None of these touch \code{drift\_moments\_jax}, \code{cg\_solve},
\code{make\_A\_apply}, or \code{m\_step\_kernel\_jax}.  The custom-VJP
shims \code{ef\_with\_jac\_grad} and \code{eff\_with\_grads} are
unchanged.

\end{document}
